\documentclass[11pt]{article}
\usepackage[margin=1.1in]{geometry}
\usepackage{amsmath, amssymb}
\usepackage{booktabs}
\usepackage{graphicx}
\usepackage[hidelinks]{hyperref}
\usepackage{caption}
\captionsetup{font=small, labelfont=bf}

\newcommand{\rhat}{\widehat{RC}}
\newcommand{\that}[1]{\widehat{\theta}_{#1}}

\title{Do Structural Models Recover Mechanisms?\\
\large Evidence from Neural Agents with Observable Internals}
\author{Briac Sockalingum\thanks{briac@berkeley.edu. Research note.
Code and full results: \url{https://github.com/briacSck/structural-interp}.
The subsidy-counterfactual experiment in Section~\ref{sec:subsidy} was
pre-registered in the repository before execution.}}
\date{July 2026 --- Working note}

\begin{document}
\maketitle

\begin{abstract}
\noindent Structural econometrics earns its keep through counterfactuals,
which are valid only if estimated parameters capture the agent's decision
\emph{mechanism} rather than merely fitting its behavior. With human data
this premise is untestable: the mechanism is never observed. We build a
laboratory where it is. Neural-network agents trained on a canonical dynamic
discrete choice environment (a two-state-variable version of Rust's (1987)
bus-engine replacement problem) are fully inspectable with
mechanistic-interpretability tools, and each agent instantiates a named
bounded-rationality primitive: rational-inattention attenuation
(Mat\v{e}jka--McKay), sparse-max attention (Gabaix), categorical cognition
(Fryer--Jackson), and salience asymmetry. Three findings. First, structural
estimation is \emph{robust} to bounded rationality that stays inside the
model class: inattentive agents' attention is absorbed into as-if
parameters, and counterfactuals --- including a pre-registered Lucas-style
subsidy experiment in which agents re-learn --- remain accurate. Second,
estimation fails \emph{silently} for cognition outside the model class:
categorical agents with 2--3\% behavioral gaps yield replacement-cost
estimates biased by up to a factor of two and policy-response predictions
wrong by a factor of two. Third, linear probes cannot tell these cases apart
(decodability is not use), but a battery of \emph{causal conformity audits}
--- the dispersion of locally patched sensitivities per state variable,
computed on in-distribution activations only --- rank-predicts
counterfactual failure (Spearman 0.87 for out-of-distribution transport,
0.80 for the subsidy experiment). Internal audits, we argue, are usefully
understood as specification tests on mechanisms, and their scope is
delimited exactly by the invariance of the cognitive primitive.
\end{abstract}

\section{Introduction}\label{sec:intro}

A structural econometrician who estimates a dynamic discrete choice (DDC)
model does not merely summarize behavior; she claims to have recovered
\emph{deep} parameters --- maintenance costs, replacement costs, attention to
economic conditions --- that remain invariant when policy changes, and that
therefore license counterfactual prediction \citep[in the sense
of][]{lucas1976, heckman2010}. The claim has a well-known analytical
vulnerability: DDC models are not nonparametrically identified
\citep{magnacthesmar2002}, identification of counterfactuals is delicate
\citep{kalouptsidi2021}, and local misspecification translates into
first-order bias \citep{bugniura2019}. What the literature has lacked is a
setting in which the vulnerability can be \emph{observed}: with human
subjects one never sees the true decision mechanism, so one can never say,
for a given estimated model, whether its parameters describe the mechanism
or merely fit the choices.

This note builds such a setting. Our agents are small neural networks
trained to implement decision rules in a solved dynamic environment. Their
behavior can be simulated at scale; their parameters --- every weight ---
are observable; and the tools of mechanistic interpretability
\citep{alainbengio2016, heimersheimnanda2024, elazar2021} allow causal
interventions on their internal representations. Crucially, we do not
construct arbitrary decision rules. Each agent instantiates a
bounded-rationality primitive with a name and a literature: rational
inattention in the logit family \citep{matejkamckay2015, steiner2017},
sparse-max attention \citep{gabaix2014}, categorical cognition
\citep{fryerjackson2008}, and asymmetric salience. By the equivalence
results of \citet{fosgerau2020}, the inattentive agents remain
\emph{inside} the parametric class the econometrician assumes, while the
categorical agents do not --- and this distinction, we will show, is exactly
what determines counterfactual validity.

The experiment has the structure of a sting operation. All agents are
behaviorally similar where the data live (ergodic choice-probability gaps of
1--3\%). We hand their simulated panels to a standard estimator (nested
fixed point maximum likelihood; a Hotz--Miller CCP estimator gives the same
picture), obtain plausible, well-converged estimates for every agent, and
then ask two questions. \emph{Ex post}: do the estimated models predict
behavior under counterfactuals --- an out-of-distribution regime shift, and
a pre-registered 50\% replacement subsidy under which agents re-learn?
\emph{Ex ante}: could an auditor with access to the agent's internals, but
only in-distribution data, have predicted which counterfactuals would fail?

Three results. \textbf{First, in-class bounded rationality is benign.} The
estimator absorbs inattention into as-if parameters --- the estimated
regime-sensitivity $\that{2}$ tracks the attention parameter across the
sparse-max and rational-inattention families almost linearly --- and the
resulting counterfactuals are correct \emph{for those agents}, including
under the subsidy experiment with re-learning. A fully regime-blind
sparse-max agent is fit with $\that{2}\approx 0$ and its policy predictions
are the most accurate in our zoo. This is \citet{fosgerau2020} made
empirical, and it is good news for structural practice that we have not seen
stated this way. \textbf{Second, out-of-class cognition fails silently.}
Categorical agents --- lookup tables over coarse state categories --- are
projected onto the parametric class: replacement-cost estimates fall to
4.6--7.0 against a truth of 10, nothing in the likelihood signals a problem,
and the predicted response to the subsidy is roughly twice the actual
response. \textbf{Third, causal internal audits see the difference where
probes cannot.} Linear probes decode the environment state from every
agent's activations ($R^2 \approx 1$) and are useless; their directions are
causally idle, reproducing the amnesic-probing lesson of \citet{elazar2021}
in an economic setting. But a \emph{conformity audit} --- the dispersion of
locally patched causal sensitivities along each state variable, gated by
insensitivity --- rank-predicts counterfactual failure at Spearman 0.87
(regime transport) and 0.80 (subsidy), using in-distribution activations
only.

We state the honest position of these results in the literature up front.
That DDC counterfactuals can fail under misspecification is known theory
\citep{magnacthesmar2002, kalouptsidi2021, bugniura2019}; that sequence
models can fit behavior with heuristics that fail to generalize is known in
machine learning \citep{vafa2024, vafa2025}; that decodability does not
imply causal use is the amnesic-probing thesis \citep{elazar2021}. Our
contribution is the synthesis: a ground-truth laboratory in which the
analytical failure becomes \emph{visible at the level of mechanism}, an
economically microfounded taxonomy of when it occurs (in-class versus
out-of-class cognition), and a demonstration that internal causal audits
constitute a workable \emph{ex-ante specification test} --- a category of
test that cannot exist for human agents but will exist for the algorithmic
agents increasingly making pricing, credit, and hiring decisions.

\section{Related work}\label{sec:lit}

\paragraph{World models and behavioral probes.} \citet{vafa2024} develop
behavioral metrics for the world model implicit in a generative model;
\citet{vafa2025} show foundation models trained on orbital trajectories
predict well yet fail to carry Newtonian mechanics to new tasks ---
``task-specific heuristics that fail to generalize,'' the phenomenon our
categorical agents exhibit. Both operate on behavior. We go inside the
network, which is what permits the \emph{ex-ante} audit. The closest
interpretability precedents are Othello-GPT \citep{li2023, nanda2023},
which established that sequence models can carry causally used world
representations, and \citet{mcgrath2022}, who probe AlphaZero for chess
concepts including value. None of these touch estimation: no econometric
procedure, no parameter recovery, no counterfactual-validity claim.

\paragraph{Identification and misspecification in DDC.}
\citet{magnacthesmar2002} establish nonparametric non-identification;
\citet{kalouptsidi2021} characterize which counterfactuals remain identified;
\citet{bugniura2019} derive the bias from local misspecification. Our
laboratory realizes these results concretely --- the categorical agents are
the misspecification, the biased $\rhat$ is the predicted consequence ---
and adds the diagnostic layer. In machine-learning dialect the same problem
is the partial identifiability of rewards in inverse reinforcement learning
\citep{skalseabate2024}.

\paragraph{Bounded rationality.} Our agent zoo is a set of named
primitives: RI-logit \citep{matejkamckay2015, steiner2017}, sparse-max
\citep{gabaix2014}, categorical cognition \citep{fryerjackson2008}. The
general equivalence of \citet{fosgerau2020} explains our central taxonomy:
attenuation-type inattention is observationally a member of the logit class,
hence absorbable into parameters; piecewise-constant categorization is not.

\paragraph{Interpretability methodology.} Linear probes
\citep{alainbengio2016} establish decodability only; amnesic probing
\citep{elazar2021} and the concept-erasure line INLP--LEACE
\citep{ravfogel2020, belrose2023} establish that decodability and use come
apart, and supply the principled erasure operators we adopt; activation
patching \citep{heimersheimnanda2024} is our causal instrument.

\section{The laboratory}\label{sec:lab}

\subsection{Environment}
The ground truth is a two-state-variable extension of \citet{rust1987}.
Mileage $x \in \{0,\dots,59\}$ accumulates stochastically; an observable
cost regime $c \in \{0,\dots,4\}$ (fuel price, labor cost, compliance
regime) follows a persistent Markov chain. Each period the agent keeps or
replaces the engine:
\[
u\big((x,c),\,\text{keep}\big) = -\theta_1\, x\Big(1 + \theta_2\,
\tfrac{c}{4}\Big), \qquad
u\big((x,c),\,\text{replace}\big) = -RC,
\]
with additive type-I extreme-value shocks, discount factor $\beta = 0.95$,
and truth $(\theta_1, \theta_2, RC) = (0.05,\, 0.8,\, 10)$. Replacement
resets mileage. The model is solved exactly by value iteration, giving the
true value function, conditional choice probabilities (CCPs), and optimal
policy.

\subsection{Agents: a microfounded zoo}
All thirteen agents share one architecture --- a two-hidden-layer MLP
mapping state features to a replacement probability --- and differ only in
the decision rule they are trained to implement:
the \emph{clone} (optimal CCPs); six \emph{categorical} agents that act on
ergodic-weighted category averages of the optimal policy over grids from
$4\times1$ to $16\times5$ \citep{fryerjackson2008}; two \emph{sparse-max}
agents perceiving $c' = m c + (1-m)\bar c$ with attention $m \in \{0, 0.5\}$
\citep{gabaix2014}; two \emph{RI} agents attenuating toward the regime-blind
policy with weight $\lambda \in \{0.25, 0.75\}$ \citep{matejkamckay2015,
fosgerau2020}; and two \emph{salience-asymmetric} agents that ignore
regimes below (above) the median. Identical architectures and inputs mean
that any difference in whether a variable is \emph{causally used} lives in
the learned weights --- precisely what an internal audit must detect.

\subsection{The econometrician}
For each agent we simulate a panel (500 buses $\times$ 200 periods,
reproducible seeds) and estimate $(\theta_1, \theta_2, RC)$ by nested
fixed point maximum likelihood, treating transition parameters as known.
(A Hotz--Miller CCP estimator exploiting replacement as a renewal action
gives near-identical estimates in the one-dimensional pilot; we report
NFXP throughout.) On the clone's data the estimator recovers
$(0.047,\, 0.85,\, 9.53)$ --- the pipeline is validated before any claim
rests on it.

\section{Results}\label{sec:results}

\subsection{Confidently wrong: fit does not imply mechanism}
\label{sec:wrong}
Table~\ref{tab:zoo} reports the zoo. Behavioral gaps on the ergodic state
distribution are 1--3.5\% --- an econometrician with a realistic panel could
not tell these agents apart --- and every estimation converges to a
plausible interior optimum. Yet the categorical agents' estimates are far
from the truth: $\rhat$ between 4.6 and 7.0 against 10, $\that{1}$
attenuated by up to 60\%. Their out-of-distribution transport fails
correspondingly: the regime-shift counterfactual (predicting behavior with
the economy pinned at the worst regime) has RMSE up to 0.218, versus 0.008
for the clone. Nothing in the fit signals this: the likelihood is smooth,
the parameters are interior, standard errors would be unremarkable.

\begin{table}[t]
\centering\small
\caption{The microfounded zoo: behavioral similarity, structural estimates,
counterfactual failure, and the pre-policy conformity audit. Truth:
$(\theta_1,\theta_2,RC) = (0.05, 0.8, 10)$. CF RMSE is the regime-shift
transport error; ``subsidy RMSE'' the policy-prediction error in the
pre-registered subsidy experiment (Section~\ref{sec:subsidy}).}
\label{tab:zoo}
\begin{tabular}{lccccccc}
\toprule
Agent & Behav.\ gap & $\that{1}$ & $\that{2}$ & $\rhat$ & CF RMSE &
Subsidy RMSE & Conformity \\
\midrule
clone        & 0.001 & 0.047 & 0.85 & 9.53  & 0.008 & 0.013 & 0.45 \\
cat$4{\times}1$  & 0.035 & 0.019 & 0.01 & 4.59  & 0.218 & 0.252 & 1.00 \\
cat$6{\times}1$  & 0.026 & 0.034 & 0.05 & 6.30  & 0.118 & 0.130 & 0.80 \\
cat$6{\times}2$  & 0.024 & 0.029 & 0.63 & 6.54  & 0.107 & 0.122 & 1.00 \\
cat$8{\times}2$  & 0.014 & 0.033 & 0.59 & 7.01  & 0.075 & 0.096 & 0.97 \\
cat$12{\times}3$ & 0.006 & 0.044 & 0.68 & 8.72  & 0.023 & 0.031 & 0.47 \\
cat$16{\times}5$ & 0.003 & 0.045 & 0.83 & 9.20  & 0.017 & 0.020 & 0.45 \\
sparse $m{=}0$   & 0.019 & 0.071 & 0.00 & 10.05 & 0.005 & 0.004 & 0.46 \\
sparse $m{=}0.5$ & 0.009 & 0.062 & 0.34 & 10.17 & 0.008 & 0.007 & 0.42 \\
RI $\lambda{=}0.25$ & 0.014 & 0.064 & 0.13 & 9.74 & 0.008 & 0.008 & 0.34 \\
RI $\lambda{=}0.75$ & 0.005 & 0.055 & 0.54 & 9.90 & 0.005 & 0.004 & 0.44 \\
asym low     & 0.012 & 0.065 & 0.28 & 9.91  & 0.026 & 0.016 & 0.70 \\
asym high    & 0.008 & 0.052 & 0.42 & 9.86  & 0.023 & 0.019 & 0.53 \\
\bottomrule
\end{tabular}
\end{table}

\subsection{In-class bounded rationality is absorbed, not punished}
\label{sec:inclass}
The striking rows of Table~\ref{tab:zoo} are the inattentive agents. The
fully regime-blind sparse-max agent is fit with $\that{2} = 0.005$ and
$\rhat = 10.05$, and has the \emph{lowest} counterfactual error in the zoo.
Across both attention families, $\that{2}$ tracks the attention parameter
($m$: $0.00 \to 0.34$; $\lambda$: $0.13 \to 0.54$): the estimator recovers
the agent's \emph{as-if} preferences, and counterfactuals computed from them
are correct \emph{for that agent} --- if the agent does not perceive the
regime, a model that says so predicts its behavior under regime shifts
perfectly. The mechanism behind this robustness is the observational
equivalence of \citet{fosgerau2020}: attenuation-type inattention stays
within the logit class, so it maps into parameters rather than into
misspecification. The practical implication for structural work is
double-edged. Reassuring: inattention of this type does not invalidate
counterfactuals. Alarming: ``uses the variable less than the modeler
assumed'' is invisible \emph{and harmless}, so audits that test variable
use --- including our own LEACE-based amnesic audit, which measures exactly
that --- are answering the wrong question. Rank correlation between the
regime-use audit and counterfactual failure across the zoo: $-0.01$.

\subsection{Probes fail; probe directions are causally idle}
\label{sec:probes}
Linear probes \citep{alainbengio2016} decode the value function and the
regime from every agent's hidden layers with $R^2 \geq 0.99$ --- including
from agents that demonstrably never use the regime. On a low-dimensional
state this is unavoidable (any monotone representation decodes anything
monotone), and it makes decodability worthless as a mechanism diagnostic.
The failure is causal as well as statistical: patching the probe direction
--- transplanting only its component between states --- moves behavior
essentially not at all (recovery of the behavioral gap $\approx 0.03$),
while patching the network's readout direction recovers 0.8--1.0 of it.
A direction can carry a concept perfectly and carry none of the decision.
This reproduces, in the smallest possible economic setting, the
amnesic-probing thesis of \citet{elazar2021} and the decodability/causality
dissociation the concept-erasure literature predicts \citep{ravfogel2020,
belrose2023}.

\subsection{Conformity audits: causal dispersion predicts failure}
\label{sec:audit}
What distinguishes the fatal agents is not \emph{which} variables they use
but \emph{how}: lookup tables are locally flat and jump at category
boundaries, while every member of the assumed model class responds smoothly.
Our audit measures this directly. For each state variable, we patch the
readout-direction component of the activation shift between adjacent states
(mileage steps $x \to x{+}2$; regime steps $c \to c{+}1$) at in-distribution
states, record the causal effect on the replacement probability, and compute
the \emph{dispersion} (coefficient of variation) of these local
sensitivities. A channel to which the agent is uniformly insensitive is
scored as conforming --- some parameter value fits it exactly, which is
precisely why regime-blindness is benign (Section~\ref{sec:inclass}); this
``insensitivity gate'' is theoretically load-bearing, not a convenience
(without it the audit misranks the regime-blind agent and the rank
correlation drops from 0.87 to 0.69). The conformity score is the worse of
the two normalized channel dispersions.

The audit uses \emph{only} in-distribution activations --- no counterfactual
data, no knowledge of the true parameters. It rank-predicts regime-shift
counterfactual failure across the thirteen agents at Spearman
\textbf{0.87} (single channels: 0.53 and 0.18). Within the categorical
family, where coarseness varies cleanly, the ordering is perfect (Spearman
1.00, $n=6$). The residual imperfection is honest and documented: the
regime-blind sparse-max agent's mileage-sensitivity profile is somewhat
uneven, which the audit reads as mild lookup-ness.

\subsection{A pre-registered Lucas experiment: the subsidy counterfactual}
\label{sec:subsidy}
Transport tests hold the agent fixed. The sharper question --- the Lucas
question --- is whether pre-policy estimates predict behavior after an
intervention to which agents \emph{adapt}. We pre-registered the design and
hypotheses in the repository before running it: a 50\% proportional
replacement subsidy ($RC\!: 10 \to 5$); adaptation under a
\emph{stable-cognitive-primitive} rule (each agent re-derives its policy in
the new environment holding its cognitive constraint fixed --- same
categories, same attention, same clamp); the econometrician predicts with
pre-policy estimates at $\rhat/2$; metrics are the policy RMSE and the error
in the predicted post-subsidy replacement rate --- the policy elasticity a
regulator actually needs.

Both hypotheses were confirmed. In-class agents pass: policy RMSE
0.004--0.013, replacement-rate errors below 0.002. Categorical agents fail,
ordered by coarseness: for the coarsest, the model predicts a post-subsidy
replacement rate of 0.134 against an actual 0.094 --- \textbf{a predicted
policy response roughly twice the actual one}. A cost-benefit analysis of
the subsidy built on that structural model would overstate its effect by
that factor while resting on an excellent in-sample fit. The pre-policy
conformity audit rank-predicts subsidy failure at Spearman \textbf{0.80}
(0.63 for the rate error).

\section{Discussion}\label{sec:discussion}

\paragraph{Audits as specification tests.} The results suggest reading
internal audits not as ``interpretability'' in the explanatory sense but as
\emph{specification tests on mechanisms}: each audit checks the conformity
of the agent's internal computation to the assumed model class along one
channel, and each counterfactual is threatened exactly by the channels it
stresses --- an operational echo of \citet{kalouptsidi2021}'s point that
identification is a property of specific counterfactuals, not of models.
For human agents such tests are impossible. For algorithmic agents ---
pricing systems, credit scorers, trading policies --- they are not, and
regulators contemplating structural models of algorithmic behavior could in
principle demand them.

\paragraph{The scope boundary.} Our subsidy experiment holds the cognitive
primitive fixed under intervention, and that assumption delimits the method.
If attention is \emph{endogenous} --- responding to stakes, as rational
inattention theory says it should \citep{mackowiak2023} --- then a subsidy
that changes stakes moves the primitive itself: as-if parameters estimated
pre-policy are no longer invariant, in-class agents fail too, and a
pre-policy audit cannot see it coming, because pre-policy internals are
genuinely conforming. This is the Lucas critique reconstituted one level
up --- for cognitive parameters rather than behavioral ones --- and it
defines precisely what conformity audits certify: invariance under
\emph{mechanism-preserving} interventions. We develop the endogenous case
in companion work.

\paragraph{Limitations.} The laboratory is deliberately small: a
two-variable state, two-layer MLPs, thirteen agents, transition parameters
treated as known, and one estimator family. These choices buy exactness ---
closed ground truth, perfect behavioral control, full internal access ---
at the price of scale, and the diagnostic's behavior on deep networks,
richer state spaces, and estimated transitions is open. The audit
normalization is cross-agent; an absolute calibration (dispersion relative
to a smooth-fit baseline) is the natural next step. We view the note as
establishing the \emph{logic} of mechanism-level validation, on the model
of Othello-GPT establishing the logic of world-model probing at toy scale.

\section{Conclusion}\label{sec:conclusion}
In a laboratory where the truth is known, structural estimation proved both
more robust and more fragile than its reputation: robust to inattentive
agents it can absorb into as-if parameters, silently wrong for categorical
agents it cannot --- with fitted likelihoods that look identical. The
difference is invisible to behavioral fit and to linear probes, and visible
to causal conformity audits computed entirely in distribution. As economic
decisions migrate to artificial agents whose internals are, for the first
time in the history of the discipline, observable, mechanism-level
specification testing becomes a real option --- and, we suspect, eventually
a standard one.

\begin{thebibliography}{99}\small

\bibitem[Alain \& Bengio(2016)]{alainbengio2016}
Alain, G., \& Bengio, Y. (2016). Understanding intermediate layers using
linear classifier probes. \emph{arXiv:1610.01644}.

\bibitem[Belrose et al.(2023)]{belrose2023}
Belrose, N., Schneider-Joseph, D., Ravfogel, S., Cotterell, R., Raff, E.,
\& Biderman, S. (2023). LEACE: Perfect linear concept erasure in closed
form. \emph{NeurIPS 2023}; arXiv:2306.03819.

\bibitem[Bugni \& Ura(2019)]{bugniura2019}
Bugni, F.~A., \& Ura, T. (2019). Inference in dynamic discrete choice
problems under local misspecification. \emph{Quantitative Economics},
10(1), 67--103.

\bibitem[Elazar et al.(2021)]{elazar2021}
Elazar, Y., Ravfogel, S., Jacovi, A., \& Goldberg, Y. (2021). Amnesic
probing: Behavioral explanation with amnesic counterfactuals.
\emph{Transactions of the ACL}, 9, 160--175.

\bibitem[Estrella \& Fuhrer(1999)]{estrellafuhrer1999}
Estrella, A., \& Fuhrer, J.~C. (1999). Are ``deep'' parameters stable? The
Lucas critique as an empirical hypothesis. \emph{FRB Boston WP} 99-4.

\bibitem[Fosgerau et al.(2020)]{fosgerau2020}
Fosgerau, M., Melo, E., de Palma, A., \& Shum, M. (2020). Discrete choice
and rational inattention: A general equivalence result.
\emph{International Economic Review}, 61(4), 1569--1589.

\bibitem[Fryer \& Jackson(2008)]{fryerjackson2008}
Fryer, R., \& Jackson, M.~O. (2008). A categorical model of cognition and
biased decision making. \emph{B.E.\ Journal of Theoretical Economics}, 8(1).

\bibitem[Fudenberg \& Strzalecki(2015)]{fudenbergstrzalecki2015}
Fudenberg, D., \& Strzalecki, T. (2015). Dynamic logit with choice
aversion. \emph{Econometrica}, 83(2), 651--691.

\bibitem[Gabaix(2014)]{gabaix2014}
Gabaix, X. (2014). A sparsity-based model of bounded rationality.
\emph{Quarterly Journal of Economics}, 129(4), 1661--1710.

\bibitem[Heckman(2010)]{heckman2010}
Heckman, J.~J. (2010). Building bridges between structural and program
evaluation approaches to evaluating policy. \emph{Journal of Economic
Literature}, 48(2), 356--398.

\bibitem[Heimersheim \& Nanda(2024)]{heimersheimnanda2024}
Heimersheim, S., \& Nanda, N. (2024). How to use and interpret activation
patching. \emph{arXiv:2404.15255}.

\bibitem[Hotz \& Miller(1993)]{hotzmiller1993}
Hotz, V.~J., \& Miller, R.~A. (1993). Conditional choice probabilities and
the estimation of dynamic models. \emph{Review of Economic Studies}, 60(3),
497--529.

\bibitem[Kalouptsidi et al.(2021)]{kalouptsidi2021}
Kalouptsidi, M., Scott, P.~T., \& Souza-Rodrigues, E. (2021).
Identification of counterfactuals in dynamic discrete choice models.
\emph{Quantitative Economics}, 12(2), 351--403.

\bibitem[Li et al.(2023)]{li2023}
Li, K., Hopkins, A.~K., Bau, D., Vi\'egas, F., Pfister, H., \& Wattenberg,
M. (2023). Emergent world representations: Exploring a sequence model
trained on a synthetic task. \emph{ICLR 2023}; arXiv:2210.13382.

\bibitem[Lucas(1976)]{lucas1976}
Lucas, R.~E. (1976). Econometric policy evaluation: A critique.
\emph{Carnegie-Rochester Conference Series on Public Policy}, 1, 19--46.

\bibitem[Ma\'ckowiak et al.(2023)]{mackowiak2023}
Ma\'ckowiak, B., Mat\v{e}jka, F., \& Wiederholt, M. (2023). Rational
inattention: A review. \emph{Journal of Economic Literature}, 61(1),
226--273.

\bibitem[Magnac \& Thesmar(2002)]{magnacthesmar2002}
Magnac, T., \& Thesmar, D. (2002). Identifying dynamic discrete decision
processes. \emph{Econometrica}, 70(2), 801--816.

\bibitem[Mat\v{e}jka \& McKay(2015)]{matejkamckay2015}
Mat\v{e}jka, F., \& McKay, A. (2015). Rational inattention to discrete
choices: A new foundation for the multinomial logit model. \emph{American
Economic Review}, 105(1), 272--298.

\bibitem[McGrath et al.(2022)]{mcgrath2022}
McGrath, T., Kapishnikov, A., Toma\v{s}ev, N., Pearce, A., Wattenberg, M.,
Hassabis, D., Kim, B., Paquet, U., \& Kramnik, V. (2022). Acquisition of
chess knowledge in AlphaZero. \emph{PNAS}, 119(47), e2206625119.

\bibitem[Nanda et al.(2023)]{nanda2023}
Nanda, N., Lee, A., \& Wattenberg, M. (2023). Emergent linear
representations in world models of self-supervised sequence models.
\emph{BlackboxNLP @ EMNLP 2023}; arXiv:2309.00941.

\bibitem[Ravfogel et al.(2020)]{ravfogel2020}
Ravfogel, S., Elazar, Y., Gonen, H., Twiton, M., \& Goldberg, Y. (2020).
Null it out: Guarding protected attributes by iterative nullspace
projection. \emph{ACL 2020}.

\bibitem[Rust(1987)]{rust1987}
Rust, J. (1987). Optimal replacement of GMC bus engines: An empirical model
of Harold Zurcher. \emph{Econometrica}, 55(5), 999--1033.

\bibitem[Skalse \& Abate(2024)]{skalseabate2024}
Skalse, J., \& Abate, A. (2024). Partial identifiability and
misspecification in inverse reinforcement learning.
\emph{arXiv:2411.15951}.

\bibitem[Steiner et al.(2017)]{steiner2017}
Steiner, J., Stewart, C., \& Mat\v{e}jka, F. (2017). Rational inattention
dynamics: Inertia and delay in decision-making. \emph{Econometrica}, 85(2),
521--553.

\bibitem[Vafa et al.(2024)]{vafa2024}
Vafa, K., Chen, J., Rambachan, A., Kleinberg, J., \& Mullainathan, S.
(2024). Evaluating the world model implicit in a generative model.
\emph{NeurIPS 2024}; arXiv:2406.03689.

\bibitem[Vafa et al.(2025)]{vafa2025}
Vafa, K., Chang, P., Rambachan, A., \& Mullainathan, S. (2025). What has a
foundation model found? Using inductive bias to probe for world models.
\emph{ICML 2025}; arXiv:2507.06952.

\end{thebibliography}

\end{document}
